\documentclass[12pt]{article}

%----------------- Packages ------------------
\usepackage[utf8]{inputenc}
\usepackage[T1]{fontenc}
\usepackage{amsmath, amssymb}
\usepackage{geometry}
\usepackage{xcolor}
\usepackage{titlesec}
\usepackage{fancyhdr}
\usepackage[breakable]{tcolorbox}
\usepackage{graphicx}
\usepackage{tikz}
\usepackage{float}
\usetikzlibrary{arrows.meta, decorations.markings}

%----------------- Page Setup -----------------
\geometry{a4paper, margin=2.5cm}
\pagestyle{fancy}
\fancyhf{}
\rhead{Final Exam — \textbf{2023}}
\lhead{Analysis IV}
\cfoot{\thepage}

%----------------- Title Styling --------------
\titleformat{\section}{\normalfont\Large\bfseries}{Exercise \thesection:}{1em}{}
\titleformat{\subsection}{\normalfont\bfseries}{Answer:}{1em}{}

%----------------- Custom Boxes ----------------
\tcbuselibrary{listingsutf8}
\newtcolorbox{answerbox}{
  colback=gray!10,
  colframe=black,
  fonttitle=\bfseries,
  title=Answer Area,
  breakable,
  before skip=10pt,
  after skip=10pt
}

%----------------- Document Start --------------
\begin{document}

%----------------- Exam Info -------------------
\begin{center}
  \Large\textbf{UNIVERSITY NAME} \\[1em]
  \large\textit{Analysis IV — Final Exam} \\[0.5em]
  \large\textit{Year: 2023} \\[2em]
\end{center}

\vspace{0.5cm}

%----------------- EXERCISE 1 ------------------
\section{}
Determine the nature (convergence/divergence) of the following numerical series:

\[
\begin{aligned}
&1) \; u_n = \frac{n^n}{n!} \\
&2) \; u_n = \log\left(1 + \frac{2}{n(n+1)}\right) \\
&3) \; u_n = \sin \left(\frac{\pi}{2} \frac{a n + 1}{\sqrt{n^2 + a n + 1}}\right)
\end{aligned}
\]

\begin{answerbox}

\end{answerbox}

%----------------- EXERCISE 2 ------------------
\section{}
For all \( n \in \mathbb{N}^* \), define the function:

\[
f_n : \mathbb{R} \longrightarrow \mathbb{R}, \quad x \mapsto \frac{n^\alpha x}{1 + n^{2} x^2}
\]

When the series \( \sum f_n(x) \) converges, denote \( f_\alpha(x) = \sum_{n=1}^{\infty} f_n(x) \).

\begin{enumerate}
    \item Determine, according to the values of \( \alpha \), the domain \( D_\alpha \) of definition of \( f_\alpha \).
    
    From now on, assume \( \alpha < 1 \).
    
    \item 
    \begin{enumerate}
        \item Show that for all \( x \in \mathbb{R}_+^* \),
        \[
        \frac{x}{1+x^2} n^{\alpha-2}\leq f_n(x) \leq \frac{1}{x} n^{\alpha-2}.
        \]
        \item Show that there exists \( A > 0 \) (depending on \( \alpha \)) such that
        \[
        \forall x \in \mathbb{R}_+^*, \quad \frac{A x}{1+x^2} \leq f_\alpha(x) \leq \frac{A}{x}.
        \]
        Deduce an equivalent of \( f_\alpha \) at \( \infty \).
    \end{enumerate}
    
    \item If \( \alpha < 0 \), show that
    \[
    \int_0^1 f_\alpha(t) \, dt = \sum_{n=1}^{\infty} \frac{1}{2} n^{\alpha-2}\log(1+n^{2}).
    \]
    
    \item Assume \( 0 \leq \alpha < 1 \).
    \begin{enumerate}
        \item Show that \( f_\alpha \) is continuous on \( \mathbb{R}^+ \).
        
        \item Show that for all \( N \in \mathbb{N}^* \)
        \[
        f_{\alpha}\left(\frac{1}{N}\right) \ge N \sum_{n=1}^{N} \frac{n^{\alpha}}{N^2 + n^2}
        \]
        Deduce the discontinuity of \( f_\alpha \) at 0.
        
        \item Does the series \( \sum f_n \) converge uniformly on \( \mathbb{R}_+ \)?
    \end{enumerate}
    
    \item (Optional) Show that \( f_\alpha \) is differentiable on \( (0, +\infty) \) and decreasing on \( [1, +\infty) \).
\end{enumerate}

\begin{answerbox}

\end{answerbox}

%----------------- EXERCISE 3 ------------------
\section{}
Consider the differential equation:
\[
4x y''(x) + 2y'(x) - y(x) = 0 \quad (E)
\]
Let \( \sum a_n x^n \) be a power series with radius of convergence \( R \) and sum \( f \).

\begin{enumerate}
    \item Justify that \( f \) is of class \( C^2 \) on \( (-R, R) \). Deduce, for \( x \in (-R, R) \), an expression of \( f'(x) \) and \( f''(x) \).
    
    \item Show that \( f \) is a solution of \( (E) \) if and only if
    \[
    a_{n+1} = \frac{a_n}{(2n+1)(2n+2)}, \quad \forall n \in \mathbb{N}.
    \]
    
    \item Express \( a_n \) in terms of \( n \) and \( a_0 \). Determine the radius of convergence of the power series \( \sum a_n x^n \).
    
    \item Expand into a power series the function defined by:
    \[
    \varphi(x) = 
    \begin{cases}
    \cosh(\sqrt{x}) & \text{if } x > 0 \\
    1 & \text{if } x = 0 \\
    \cos(\sqrt{-x}) & \text{if } x < 0
    \end{cases}
    \]
\end{enumerate}

\begin{answerbox}

\end{answerbox}

%----------------- END ------------------
\end{document}